% !TeX root = ../../../main.tex
\begin{sloppypar}
\subsubsection{Related work}
Numerous critical care and sepsis datasets are available \mbox{(semi-)publicly}, \emph{e.g.}, \cite{thoral2021sharing,pollard2018eicu}, including several with de-identified clinical discharge summaries: MIMIC-III \citep{ch06:johnson2016mimic}
%,goldberger2000physiobank}
and MIMIC-IV \citep{ch06:johnson2023mimic}. 
Sepsis phenotyping (including Sepsis-3) has been conducted using a variety of methods on structured data \citep{johnson2018comparative}, although structured data is often highly missing, leading to suboptimal phenotyping \citep{seinen2025using}.
Additionally, the timing of sepsis phenotyping is crucial, in that delays in sepsis prediction can nullify predictiveness and lead to burdensome and unhelpful analytic models \citep{kamran2024evaluation}.  
\end{sloppypar}

Several works have used the i2b2 (a competition subset of MIMIC-II/III) to construct timelines of clinical concepts \citep{leeuwenberg2020towards,ch06:frattallone2024using}.  
In the case of \cite{leeuwenberg2020towards}, the event spans are already given, and focus on an excerpt of the full discharge summary, so we focus our direct comparison on that of \cite{ch06:frattallone2024using}.
% Frattallone et al.  
In contrast to these works, we adopt an annotation process focused on temporal assignment to clinical findings rather than i2b2 clinical concepts \citep{ch06:uzuner20112010}, which allows for greater specificity of the event in the absence of additional context.  
Additionally, as compared to \cite{ch06:frattallone2024using},
% Frattallone et al.
we focus on using text alone for clinical time series construction, which is necessary for the PubMed Open Access Subset since no complementary data source is available.

Broadly in the medical domain, LLMs are showing promise in medical tasks such as summarization \citep{van2024adapted}.
% ,aali2024benchmark}. 
In fact, some findings suggest that 
medical fine-tuning rarely leads to improved performance over the corresponding foundation models \citep{jeong-etal-2024-medical}. 
In our analysis, we primarily use foundation LLMs and show that their scale and instruction-tuning are central to the performance of the model.  

